% ============================================================================
% Fichier  : partie4/P04_C20_exercer.tex
% Titre    : Exercer et Tester les Plans
% Auteur   : MINKA MI NGUIDJOÏ Thierry Emmanuel
% Version  : 1.0 -- novembre 2026
% ============================================================================
% Sources : ISO 22301:2019 (8.5) ; NIST SP 800-84 (2006). Programme
% d'exercices de LiZbiZ et bascule de la BANQUE SANAGA : cas fictifs.
% ============================================================================

\chapter{Exercer et Tester les Plans}
\label{chap:exercer}
\chaptermark{Exercer et tester les plans}

\epigraph{%
    « Un exercice qui réussit tout n'a rien appris à personne. »%
}{Principe du cours}

\niveauChapitre{Intermédiaire}{Chapitres~\ref{chap:pca}
et~\ref{chap:pra}.}

\begin{boxobjectifs}
\begin{itemize}
    \item Distinguer revue, exercice et test, et leurs niveaux de réalisme.
    \item Construire un programme d'exercices sur plusieurs années.
    \item Concevoir un exercice~: objectifs mesurables, scénario, injections,
          observation.
    \item Évaluer un exercice et en tirer un plan d'actions.
    \item Limiter les risques d'un test sur des systèmes réels.
\end{itemize}
\end{boxobjectifs}

\bigskip

% ============================================================================
\section*{Question directrice : comment savoir que le plan de continuité fonctionne, avant la prochaine crise~?}
% ============================================================================

En juin, LiZbiZ dispose d'un plan de continuité et d'un plan de reprise.
Personne ne les a jamais utilisés. En août, l'inondation de Bonabéri a
prouvé le premier pour les expéditions de Douala~; mais la facturation,
les autres sites et la cellule de crise n'ont pas été éprouvés. Il faut un
programme d'exercices, avant que les événements ne s'en chargent.

\medskip
\textbf{Plan du chapitre.}
\begin{enumerate}
    \item Revue, exercice, test (section~\ref{sec:c20-types}).
    \item Le programme d'exercices (section~\ref{sec:c20-programme}).
    \item Concevoir un exercice (section~\ref{sec:c20-concevoir}).
    \item Évaluer (section~\ref{sec:c20-evaluer}).
    \item Tester sans nuire (section~\ref{sec:c20-nuire}).
\end{enumerate}

\separateur

% ============================================================================
\section{Revue, Exercice, Test}
\label{sec:c20-types}
% ============================================================================

L'ISO~22301 exige un programme d'exercices et de tests, fondé sur des
scénarios, qui valide dans le temps l'efficacité des stratégies et des
plans~\cite{iso22301_2019}. Le NIST SP~800-84 décrit les programmes de
tests, de formation et d'exercices~\cite{nist80084}.

\begin{tableaucours}{L{3.4cm}L{5.0cm}L{1.8cm}L{2.0cm}}{%
    \tableauHeader{Forme} &
    \tableauHeader{Principe} &
    \tableauHeader{Réalisme} &
    \tableauHeader{Coût}
}
Revue documentaire & Relire le plan, vérifier contacts et procédures &
Faible & Faible \\
Exercice sur table & Les acteurs discutent un scénario autour d'une table
& Moyen & Faible \\
Exercice fonctionnel & Les acteurs agissent dans leurs rôles, sans toucher
la production & Élevé & Moyen \\
Test technique & Restaurer, basculer, mesurer sur des systèmes réels ou
copiés & Élevé & Moyen à élevé \\
Exercice grandeur réelle & Toute l'organisation, systèmes compris &
Très élevé & Élevé \\
\end{tableaucours}
\vspace{-0.8em}
\begin{center}{\small\itshape Tableau~20.1. Les formes d'exercice et de test}\end{center}

\begin{boxdef}
\termecle{Exercice}~: activité qui entraîne des personnes et évalue des
plans ou des capacités à partir d'un scénario. \termecle{Test}~: exercice
dont on attend un résultat mesurable, réussi ou échoué (une restauration
en moins de 48~heures).
\end{boxdef}

% ============================================================================
\section{Le Programme d'Exercices}
\label{sec:c20-programme}
% ============================================================================

Un programme progresse du simple au complexe, couvre toutes les activités
critiques sur un cycle, et varie les scénarios.

\begin{tableaucours}{L{2.2cm}L{4.4cm}L{3.4cm}L{2.2cm}}{%
    \tableauHeader{Période} &
    \tableauHeader{Exercice} &
    \tableauHeader{Plan éprouvé} &
    \tableauHeader{Forme}
}
Sept. 2026 & Relecture et appel des contacts & Tous & Revue \\
Oct. 2026 & Prise en main de la cellule de crise & Gestion de crise & Sur
table \\
Nov. 2026 & Reprise de la facturation à Yaoundé & Reprise informatique &
Test technique \\
Fév. 2027 & Rançongiciel, cellule de crise & Gestion de crise & Sur table \\
Mai 2027 & Coupure de Garoua quatre jours & Continuité du Nord &
Fonctionnel \\
Sept. 2027 & Perte de Douala, bascule & Tous & Grandeur réelle partielle \\
\end{tableaucours}
\vspace{-0.8em}
\begin{center}{\small\itshape Tableau~20.2. Programme d'exercices 2026-2027 de LiZbiZ}\end{center}

% ============================================================================
\section{Concevoir un Exercice}
\label{sec:c20-concevoir}
% ============================================================================

\begin{boxprocedure}
\textbf{Concevoir un exercice}
\begin{enumerate}
    \item \textbf{Objectifs mesurables}~: « la facturation reprend en moins
          de 48~heures », « le DG est prévenu en moins d'une heure ».
    \item \textbf{Périmètre et participants}~: qui joue, qui observe, qui
          anime.
    \item \textbf{Scénario} réaliste, tiré du registre des risques.
    \item \textbf{Injections}~: messages qui font évoluer la situation à
          heures fixes (un journaliste appelle, la sauvegarde est corrompue).
    \item \textbf{Grille d'observation}~: pour chaque objectif, ce que
          l'observateur note et quand.
    \item \textbf{Règles de sécurité}~: mot d'arrêt, communication
          « EXERCICE » sur chaque message.
\end{enumerate}
\end{boxprocedure}

\begin{boxalert}
\textbf{Violation typique.} Un scénario connu à l'avance des joueurs, sans
surprise.

\smallskip
\textbf{Règle~:} les joueurs connaissent la date et le cadre, pas le
scénario. Au moins une injection doit mettre le plan en défaut (un
suppléant absent, un fournisseur injoignable)~: c'est là qu'on apprend.
\end{boxalert}

% ============================================================================
\section{Évaluer}
\label{sec:c20-evaluer}
% ============================================================================

\begin{boxscenario}{Opération LIZBIZ. Le test de reprise de novembre}
\textbf{Objectifs.} Facturation reprise en moins de 48~heures~; données
perdues inférieures à 4~heures~; contrôle de dix factures par la DAF.

\medskip
\textbf{Constatations.} Reprise en 19~heures~; perte de 5~heures de
données~; huit factures sur dix conformes, deux avec un taux de taxe
erroné (paramètre non restauré).

\medskip
\textbf{Analyse} (confiance élevée). Le RTO est tenu~; le RPO ne l'est pas
(Chapitre~\ref{chap:pra})~; la restauration des paramètres n'était pas
dans la procédure. Le test a fait apparaître deux écarts que rien d'autre
n'aurait révélés.
\end{boxscenario}

\begin{tableaucours}{L{4.4cm}L{2.2cm}L{5.6cm}}{%
    \tableauHeader{Objectif} &
    \tableauHeader{Résultat} &
    \tableauHeader{Action}
}
Reprise en moins de 48~h & Atteint (19~h) & Aucune \\
Perte inférieure à 4~h & Non atteint (5~h) & Journalisation des
transactions \\
Dix factures conformes & Non atteint (8/10) & Ajouter les paramètres à la
procédure \\
\end{tableaucours}
\vspace{-0.8em}
\begin{center}{\small\itshape Tableau~20.3. Évaluation du test de novembre}\end{center}

% ============================================================================
\section{Tester sans Nuire}
\label{sec:c20-nuire}
% ============================================================================

Un test sur des systèmes réels peut lui-même provoquer un incident. On le
prépare comme un changement~: fenêtre en heures creuses, sauvegarde
préalable, critère d'arrêt, retour arrière prévu et répété, information des
utilisateurs et des clients concernés.

\begin{boxobjet}{R-W1 · Bascule annuelle de la BANQUE SANAGA}
\textbf{Scénario testé~:} perte du site de Douala~; bascule du système
central sur le site de secours. \textbf{Fenêtre~:} un samedi de 22~h à
6~h. \textbf{Critère d'arrêt~:} toute anomalie sur un virement réel.
\textbf{Résultat attendu~:} bascule en moins de 2~heures, aucun écart de
solde. \textbf{Propriétaire~:} M.~EYENGA. \textbf{Indicateur~:} durée
mesurée de la bascule.
\end{boxobjet}

\begin{boxmetier}{Responsable de la continuité}
\textbf{Mesurer, pas raconter.} Le responsable fait chronométrer chaque
étape par un observateur et consigne les heures. Le rapport compare les
durées mesurées aux objectifs, étape par étape.
\end{boxmetier}

\separateur

\begin{boxresume}
\textbf{Ce qu'il faut retenir de ce chapitre~:}
\begin{itemize}
    \item \textbf{Revue, sur table, fonctionnel, technique, grandeur
          réelle}~: un réalisme et un coût croissants.
    \item \textbf{Le programme} progresse, couvre les activités critiques et
          varie les scénarios.
    \item \textbf{Un exercice} a des objectifs mesurables, des injections et
          des observateurs.
    \item \textbf{L'évaluation} compare objectif et résultat et produit des
          actions.
    \item \textbf{Un test réel} se prépare comme un changement, avec un
          critère d'arrêt.
\end{itemize}
\end{boxresume}

\begin{boxinfo}
\textbf{Points de contrôle de l'élève.} Je sais~: distinguer les formes
d'exercice~; construire un programme~; rédiger des objectifs mesurables et
des injections~; évaluer un exercice~; préparer un test réel sans risque
excessif.
\end{boxinfo}

\bigskip

\begin{boxexo}[title={\ding{45}\quad Exercice~20.1. Quelle forme~? \niveaubase}]
Classez~: appeler tous les contacts du plan~; restaurer la base de
facturation sur une machine de test~; réunir la direction autour d'un
scénario d'inondation.
\end{boxexo}

\begin{boxexo}[title={\ding{45}\quad Exercice~20.2. Mesurable \niveaubase}]
Réécrivez de manière mesurable~: « vérifier que la cellule de crise sait
se réunir ».
\end{boxexo}

\begin{boxexo}[title={\ding{45}\quad Exercice~20.3. Le test de novembre \niveaubase}]
Quels objectifs du test de novembre sont atteints, lesquels ne le sont
pas~?
\end{boxexo}

\begin{boxexo}[title={\ding{45}\quad Exercice~20.4. Injections \niveaumoyen}]
Proposez cinq injections pour l'exercice sur table « rançongiciel » de
février 2027, avec leur heure.
\end{boxexo}

\begin{boxexo}[title={\ding{45}\quad Exercice~20.5. Observation \niveaumoyen}]
Construisez la grille d'observation de l'exercice « coupure de Garoua
quatre jours ».
\end{boxexo}

\begin{boxexo}[title={\ding{45}\quad Exercice~20.6. Les paramètres \niveaumoyen}]
Deux factures sur dix avaient un taux erroné. Quelle cause, et quelle
correction à la procédure~?
\end{boxexo}

\begin{boxexo}[title={\ding{45}\quad Exercice~20.7. Critère d'arrêt \niveaumoyen}]
Pourquoi un critère d'arrêt est-il indispensable à la bascule de la
BANQUE SANAGA~? Que se passe-t-il s'il est atteint~?
\end{boxexo}

\begin{boxexo}[title={\ding{45}\quad Exercice~20.8. Programme \niveaumoyen}]
Le programme du Tableau~20.2 couvre-t-il toutes les activités critiques de
LiZbiZ~? Que manque-t-il~?
\end{boxexo}

\begin{boxexo}[title={\ding{45}\quad Exercice~20.9. Un dossier d'exercice \niveauexpert}]
Rédigez le dossier complet de l'exercice fonctionnel de mai 2027~:
objectifs, scénario, injections, participants, observation, sécurité.
\end{boxexo}

\begin{boxexo}[title={\ding{45}\quad Exercice~20.10. Rapport d'exercice \niveauexpert}]
Rédigez le rapport du test de novembre~: cadre, objectifs, résultats
chronométrés, écarts, actions.
\end{boxexo}

\bigskip

\begin{boxinfo}
\textbf{Pour aller plus loin.}
\begin{itemize}
    \item ISO~22301, programme d'exercices~\cite{iso22301_2019}.
    \item NIST SP~800-84~\cite{nist80084}.
\end{itemize}
\end{boxinfo}

\begin{boxinfo}
\textbf{Transition.} Quand l'incident déborde les plans, l'organisation
entre en crise. La Partie~V commence par définir la crise
(Chapitre~\ref{chap:crise}).
\end{boxinfo}
